\documentclass[../infufrgs/ppgc-tese.tex]{subfiles}
\begin{document}
\section*{Notation and model setup}
This chapter establishes the minimum notation and model setup used throughout the thesis to discuss sparse auto-encoders (SAEs), cross-layer transcoders, and logic-guided steering. We avoid a tutorial treatment of Transformers and focus instead on the representational objects we measure and manipulate.

We denote a tokenized input sequence by $x_{1:T}$ and the model dimension by $d$. A Transformer with $L$ layers produces a residual stream $h_\ell \in R^{T\times d}$ after each layer $\ell\in\{1,\dots,L\}$. Each layer applies attention and an MLP that write updates into the residual stream; we do not require head-level detail beyond noting that attention heads operate position-wise and share the residual pathway.

Unless otherwise stated, we collect activations as follows: $h_\ell$ denotes the residual stream immediately after layer $\ell$; optional probes may target head outputs or the MLP pre-/post-nonlinearity, but all core analyses use $h_\ell$. SAEs operate on $h_\ell$ to produce a sparse code $s_\ell = E_\ell(h_\ell)$ and a reconstruction $\hat h_\ell = D_\ell(s_\ell)$. We evaluate reconstruction error, sparsity of $s_\ell$, and stability across seeds/epochs.

To compare or transfer representations across depth, we define a cross-layer transcoder $T_{\ell\to k}: s_\ell \mapsto \hat s_k$, which aims to align layer-$\ell$ features with the feature basis learned at layer $k$. Logic validators provide soft scores for predicates/rules evaluated on $s_\ell$ or $\hat s_k$. Steering refers to using these scores to influence generation via logit adjustment, constrained decoding, or reranking.

\subsection*{Implications for this thesis.} The symbols $h_\ell$, $s_\ell$, and $T_{\ell\to k}$ are the canonical objects for measurement and control. All metrics (reconstruction, sparsity, stability, alignment, constraint satisfaction) are reported with respect to these objects. The steering surfaces considered later act on logits using signals derived from validators applied to $s_\ell$ or $\hat s_k$.

\section{Learning phenomena}
We briefly position the learning phenomena relevant to this work. Generalization denotes performance on unseen data; overfitting occurs when training performance improves while validation/test performance degrades; underfitting is failure to fit even the training data. In practice, we monitor paired training/validation curves and compare gaps over time to assess whether improvements reflect memorization or genuine structure acquisition.

Modern deep models often benefit from implicit regularization (e.g., optimization bias and architectural priors), which can reconcile high capacity with good generalization. Training dynamics can exhibit phases: rapid memorization, representation re-organization, and late convergence. Of particular interest is late generalization or ``grokking'': after the training loss saturates near zero, the test performance remains low for a long period and then sharply improves, stabilizing near training performance.

For our purposes, we treat grokking as a measurable delay between the end of the memorization phase and the onset of generalization, observable on controlled synthetic tasks. We track this delay alongside mechanistic signals such as the stability of sparse features $s_\ell$ and their reconstruction quality. We distinguish grokking from double descent by its temporal pattern (delayed improvement at near-zero training loss) and its association with representation re-organization rather than model capacity sweeps.

We use these phenomena in two ways: as motivation to collect activations at multiple training phases and as a source of experimental conditions that stress representation alignment across layers.

\subsection*{Implications for this thesis.} Chapter~\ref{chpt:learning_phenomena} details datasets and measurements for late generalization. Here we commit to (i) reporting stability/drift of $s_\ell$ across epochs/seeds; (ii) selecting training windows that include pre-/post-generalization phases; and (iii) using synthetic settings to elicit representation changes that test the robustness of SAEs and transcoders.

\section{Neuro-symbolic AI}
Neuro-symbolic AI spans a design space ranging from symbolic-first systems (logic programs calling learned perception modules) to neural-first systems regularized or evaluated by symbolic structure. Following Kautz's taxonomy, axes include the primary substrate (neural or symbolic), the coupling strength (tight vs.~loose), and the directionality of information flow (constraints, explanations, or supervision).

This thesis positions itself as neural-first with loose but differentiable coupling: we learn neural representations (via SAEs) and apply soft logic validators that map activations to truth scores in $[0,1]$. A validator encodes a predicate or rule as a differentiable function of features (e.g., $p(s_\ell)\in[0,1]$), aggregates them with t-norms to express conjunction/disjunction, and produces a satisfaction score or penalty. Logic Tensor Networks (LTN) are an instance of such soft logic, supporting predicate grounding over tensors and differentiable fuzzy operators.

Validators serve two roles: as training-time regularizers to favor representations that satisfy desired constraints and as evaluation-time metrics to quantify constraint satisfaction and coverage. In the cross-layer setting, validators can be shared across layers to test whether semantics are preserved by $T_{\ell\to k}$.

\subsection{Implications for this thesis.} We adopt soft logic validators (LTN-style) as both losses and metrics. They provide portable semantic checks across layers and yield signals for steering. We will report satisfaction rates and penalty magnitudes alongside reconstruction/alignment metrics to assess whether aligned features remain semantically meaningful.

\section*{Interpretability}
Interpretability goals vary along two axes: local vs.~global and mechanistic vs.~post-hoc. Local methods explain particular inputs or activations; global methods summarize model behavior or circuits. Mechanistic approaches aim to uncover causal or algorithmic structure (e.g., feature-level circuits), while post-hoc methods provide approximate explanations without claiming faithful mechanism.

Polysemanticity---multiple unrelated features entangled within a single dimension of the residual stream---motivates SAEs: by promoting sparsity and overcomplete bases, SAEs aim to separate conflated concepts into more interpretable features. However, SAEs must be evaluated for fidelity (reconstruction), usefulness (downstream control), and stability (consistency across seeds and epochs).

In this thesis, we use SAEs to obtain sparse features $s_\ell$ per layer, transcoders to align features across layers, and validators to attach semantic meaning to features. This triad supports both mechanistic inspection (feature alignment and validator transfer) and practical steering.

\subsection*{Implications for this thesis.} We prioritize mechanistic interpretability with practical utility: features should reconstruct well, be sparse and stable, align across layers, and satisfy reusable validators. These criteria drive the metrics and ablations used later.

\subsection{Sparse Auto-encoders}
A sparse auto-encoder on layer $\ell$ learns an encoder $E_\ell: h_\ell\mapsto s_\ell$ and a decoder $D_\ell: s_\ell\mapsto \hat h_\ell$ that reconstruct the residual stream from a sparse code. Sparsity encourages features that fire selectively, which can reduce polysemanticity and improve interpretability. Training typically minimizes reconstruction error with a sparsity-promoting penalty or target activation rate.

We report three primary metrics: (i) sparsity of $s_\ell$ (e.g., average activation rate), (ii) reconstruction error $\|h_\ell-\hat h_\ell\|$, and (iii) stability across seeds and epochs (feature correspondence and drift). Practical considerations include normalization of $h_\ell$, preventing dead features, and calibrating sparsity targets to balance fidelity and disentanglement.

\subsection*{Implications for this thesis.} SAE outputs $s_\ell$ constitute the feature space for alignment and validation. Stability protocols (cross-seed/epoch matching) and sparsity/reconstruction trade-offs are part of all SAE reports and ablations.

\subsection{Cross-layer Transcoder}
As depth composes functions, semantically related signals can appear at multiple layers in transformed bases. We introduce a cross-layer transcoder $T_{\ell\to k}$ to map $s_\ell$ into the basis of layer $k$, yielding $\hat s_k$. Training objectives include alignment to $s_k$ (when available), semantic consistency under validators, and regularization to prevent degenerate solutions.

Evaluation covers prediction quality of $\hat s_k$ against $s_k$, transfer of validator satisfaction across layers, and stability across seeds. Care is taken to avoid the transcoder inventing semantics not present in $s_\ell$ or $s_k$ (checked via reconstruction and validator consistency tests).

\subsection*{Implications for this thesis.} $T_{\ell\to k}$ enables cross-layer analysis and control. Success is defined by high alignment with $s_k$, stable mappings, and preserved validator semantics, which we quantify in Chapter~\ref{chpt:circuitry} and later.

\subsection{LTN Cross-layer Transcoder}
The LTN Cross-layer Transcoder combines layerwise SAEs, cross-layer mapping, and soft logic validators. SAEs produce sparse features per layer; $T_{\ell\to k}$ aligns these features; validators score whether features satisfy desired relations. These scores serve two functions: regularizing the transcoder to prefer semantically consistent alignments and producing signals for steering.

At a high level, the method aims to improve interpretability that is useful for control: aligned sparse features that reconstruct well, transfer across depth, and respect declarative constraints. Detailed objectives, training schedules, and steering surfaces are specified in Chapter~4; here we motivate the combination and its expected benefits.

\subsection*{Implications for this thesis.} This section bridges Background to Methodology. We will (i) train SAEs per layer and report sparsity/reconstruction/stability, (ii) train $T_{\ell\to k}$ and report alignment and semantic transfer, and (iii) use LTN validators both as regularizers and evaluation metrics, ultimately yielding signals for steering.

\end{document}
